\chapter{Complexity — Meyer Verifier Glue}
%%% generated by the blueprint pipeline from TCSlib.Complexity.CircuitComplexity.MeyerVerifierGlue
\section{Overview}

Declarations of the Lean module \texttt{TCSlib.Complexity.CircuitComplexity.MeyerVerifierGlue}.

\section{Declarations}

\begin{definition}[Nested first-component extraction]
\lean{Complexity.Meyer.vx}
\label{Complexity.Meyer.vx}
\leanok
\uses{Complexity.pairFstD}
For a bit list $z$, define $v_x(z)$ by extracting the first component of the first
component of $z$ using the pair decoder.
\end{definition}

\begin{definition}[Marked second-component extraction]
\lean{Complexity.Meyer.vD}
\label{Complexity.Meyer.vD}
\leanok
\uses{Complexity.KarpLipton.dropMarker, Complexity.pairFstD, Complexity.pairSndD}
For a bit list $z$, define $v_D(z)$ by taking the second component of the first
component of $z$ using the pair decoder, then removing its marker.
\end{definition}

\begin{definition}[Nested extraction from a second component]
\lean{Complexity.Meyer.vt}
\label{Complexity.Meyer.vt}
\leanok
\uses{Complexity.pairFstD, Complexity.pairSndD}
For a bit list $z$, define $v_t(z)$ by taking the first component of the first component
of the second component of $z$, using the pair decoder.
\end{definition}

\begin{definition}[Second component of a nested pair]
\lean{Complexity.Meyer.va}
\label{Complexity.Meyer.va}
\leanok
\uses{Complexity.pairFstD, Complexity.pairSndD}
For a bit list $z$, define $v_a(z)$ by taking the second component of the first
component of the second component of $z$, using the pair decoder.
\end{definition}

\begin{definition}[Nested second-component extraction]
\lean{Complexity.Meyer.vj}
\label{Complexity.Meyer.vj}
\leanok
\uses{Complexity.pairFstD, Complexity.pairSndD}
For a bit list $z$, define $v_j(z)$ by taking the first component of the second
component of the second component of $z$, using the pair decoder.
\end{definition}

\begin{definition}[Repeated second-component extraction]
\lean{Complexity.Meyer.vp}
\label{Complexity.Meyer.vp}
\leanok
\uses{Complexity.pairSndD}
For a bit list $z$, define $v_p(z)$ by taking the second component three times in
succession using the pair decoder.
\end{definition}

\begin{definition}[Nested encoding of six bit lists]
\lean{Complexity.Meyer.vin}
\label{Complexity.Meyer.vin}
\leanok
\uses{Turing.pairEncode}
For bit lists $x,U,t,a,j,p$, define their encoding by pairing the encodings of $x$ and
$U$ with the encoding of the pair consisting of the encoding of $t$ and $a$ and the
encoding of $j$ and $p$, namely
$\operatorname{pairEncode}(\operatorname{pairEncode}(x,U),\operatorname{pairEncode}(\operatorname{pairEncode}(t,a),\operatorname{pairEncode}(j,p)))$.
\end{definition}

\begin{theorem}[Recovery of the First Bit String]
\lean{Complexity.Meyer.vx_vin}
\label{Complexity.Meyer.vx_vin}
\leanok
\uses{Complexity.Meyer.vin, Complexity.Meyer.vx}
\difficulty{1}
Let $x,U,t,a,j,p$ be finite bit strings, and let $\langle u,v\rangle$ denote the
encoding formed by doubling each bit of $u$, appending the separator $01$, and then
appending $v$. If $\pi_1$ is the total first projection that returns the first component
of a pair encoding and returns the empty string otherwise, then applying $\pi_1$ twice
to $\langle\langle x,U\rangle,\langle\langle t,a\rangle,\langle j,p\rangle\rangle$ gives
$x$.
\end{theorem}

\begin{theorem}[Decoding the Nested Input]
\lean{Complexity.Meyer.vD_vin}
\label{Complexity.Meyer.vD_vin}
\leanok
\uses{Complexity.KarpLipton.dropMarker, Complexity.Meyer.vD, Complexity.Meyer.vin}
\difficulty{1}
For arbitrary bit strings $x,U,t,a,j,p$, form the nested encoding $\langle\langle
x,U\rangle,\langle\langle t,a\rangle,\langle j,p\rangle\rangle$. Applying the total
second projection, then the total first projection, then the marker-dropping transducer
to this encoding yields the result of applying that transducer to $U$.
\end{theorem}

\begin{theorem}[Recovering the Time String]
\lean{Complexity.Meyer.vt_vin}
\label{Complexity.Meyer.vt_vin}
\leanok
\uses{Complexity.Meyer.vin, Complexity.Meyer.vt}
\difficulty{1}
For arbitrary bit strings $x,U,t,a,j,p$, form the nested encoding $\langle\langle
x,U\rangle,\langle\langle t,a\rangle,\langle j,p\rangle\rangle$. Applying the total
second projection, then the total first projection twice, to this encoding recovers $t$.
\end{theorem}

\begin{theorem}[Recovering the Auxiliary String]
\lean{Complexity.Meyer.va_vin}
\label{Complexity.Meyer.va_vin}
\leanok
\uses{Complexity.Meyer.va, Complexity.Meyer.vin}
\difficulty{1}
For arbitrary bit strings $x,U,t,a,j,p$, form the nested encoding $\langle\langle
x,U\rangle,\langle\langle t,a\rangle,\langle j,p\rangle\rangle$. Applying the total
second projection, then the total first projection, then the total second projection to
this encoding recovers $a$.
\end{theorem}

\begin{theorem}[Recovering the Index String]
\lean{Complexity.Meyer.vj_vin}
\label{Complexity.Meyer.vj_vin}
\leanok
\uses{Complexity.Meyer.vin, Complexity.Meyer.vj}
\difficulty{1}
For arbitrary bit strings $x,U,t,a,j,p$, form the nested encoding $\langle\langle
x,U\rangle,\langle\langle t,a\rangle,\langle j,p\rangle\rangle$. Applying the total
second projection twice and then the total first projection to this encoding recovers
$j$.
\end{theorem}

\begin{theorem}[Recovery of the final string]
\lean{Complexity.Meyer.vp_vin}
\label{Complexity.Meyer.vp_vin}
\leanok
\uses{Complexity.Meyer.vin, Complexity.Meyer.vp}
\difficulty{1}
For any bit strings $x,U,t,a,j,p$, let $\langle r,s\rangle$ denote the encoding formed
by doubling each bit of $r$, appending the separator $01$, and then appending $s$. The
total second projection, which returns the second component of a parsed pair and the
empty string when parsing fails, applied three times to $\langle\langle
x,U\rangle,\langle\langle t,a\rangle,\langle j,p\rangle\rangle$ returns $p$.
\end{theorem}

\begin{theorem}[Polynomial-Time Computability of the Nested Projection]
\lean{Complexity.Meyer.pt_vx}
\label{Complexity.Meyer.pt_vx}
\leanok
\uses{Complexity.Meyer.vx, Complexity.PolyTimeComputable, Complexity.PolyTimeComputable.comp, Complexity.polyTimeComputable_pairFstD}
\difficulty{1}
The function on bit strings that applies the total first projection of the pair encoding
twice is polynomial-time computable.
\end{theorem}

\begin{theorem}[Polynomial-time computability of marker removal]
\lean{Complexity.Meyer.pt_vD}
\label{Complexity.Meyer.pt_vD}
\leanok
\uses{Complexity.KarpLipton.polyTimeComputable_dropMarker, Complexity.Meyer.vD, Complexity.PolyTimeComputable, Complexity.PolyTimeComputable.comp, Complexity.polyTimeComputable_pairFstD, Complexity.polyTimeComputable_pairSndD}
\difficulty{1}
For a bit string $z$, take the second component of the total first projection of $z$,
then apply the total second projection to that result; failed pair parses yield the
empty word. Finally, process the resulting word from left to right, initially outside a
marked region: a zero switches the process into the marked region and emits nothing, a
one outside the marked region emits nothing, and every bit inside the marked region is
emitted unchanged. The function so defined is polynomial-time computable.
\end{theorem}

\begin{theorem}[Polynomial-time computability of a nested projection]
\lean{Complexity.Meyer.pt_vt}
\label{Complexity.Meyer.pt_vt}
\leanok
\uses{Complexity.Meyer.vt, Complexity.PolyTimeComputable, Complexity.PolyTimeComputable.comp, Complexity.polyTimeComputable_pairFstD, Complexity.polyTimeComputable_pairSndD}
\difficulty{1}
For a bit string $z$, take its total second projection, then the total first projection
of that result, and then the total first projection once more; whenever a pair parse
fails, the corresponding projection returns the empty word. The resulting function is
polynomial-time computable.
\end{theorem}

\begin{theorem}[Polynomial-time computability of a nested projection]
\lean{Complexity.Meyer.pt_va}
\label{Complexity.Meyer.pt_va}
\leanok
\uses{Complexity.Meyer.va, Complexity.PolyTimeComputable, Complexity.PolyTimeComputable.comp, Complexity.polyTimeComputable_pairFstD, Complexity.polyTimeComputable_pairSndD}
\difficulty{1}
For a bit string $z$, take its total second projection, then the total first projection
of that result, and finally the total second projection once more; whenever a pair parse
fails, the corresponding projection returns the empty word. The resulting function is
polynomial-time computable.
\end{theorem}

\begin{theorem}[Polynomial-time computability of a nested projection]
\lean{Complexity.Meyer.pt_vj}
\label{Complexity.Meyer.pt_vj}
\leanok
\uses{Complexity.Meyer.vj, Complexity.PolyTimeComputable, Complexity.PolyTimeComputable.comp, Complexity.polyTimeComputable_pairFstD, Complexity.polyTimeComputable_pairSndD}
\difficulty{1}
For a bit string $z$, take its total second projection twice and then take the total
first projection of the result; whenever a pair parse fails, the corresponding
projection returns the empty word. The resulting function is polynomial-time computable.
\end{theorem}

\begin{theorem}[Polynomial-time computability of a nested projection]
\lean{Complexity.Meyer.pt_vp}
\label{Complexity.Meyer.pt_vp}
\leanok
\uses{Complexity.Meyer.vp, Complexity.PolyTimeComputable, Complexity.PolyTimeComputable.comp, Complexity.polyTimeComputable_pairSndD}
\difficulty{1}
For a bit string $z$, take its total second projection three times, returning the empty
word at any stage where the pair parse fails. The resulting function is polynomial-time
computable.
\end{theorem}

\begin{theorem}[Polynomial-time computability of fixed-width increment]
\lean{Complexity.Meyer.pt_inc}
\label{Complexity.Meyer.pt_inc}
\leanok
\uses{Complexity.PolyTimeComputable, Complexity.polyTimeComputable_of_linear, Turing.FinTM.computesFunInTime_incFixed, Turing.incFixed}
\difficulty{3}
The function sending a Boolean word to its fixed-width binary increment, with overflow
replaced by the empty word, is polynomial-time computable.
\end{theorem}

\begin{theorem}[Polynomial-time computability of binary length]
\lean{Complexity.Meyer.pt_lengthBits}
\label{Complexity.Meyer.pt_lengthBits}
\leanok
\uses{Complexity.PolyTimeComputable, Complexity.polyTimeComputable_of_linear, Turing.FinTM.computesFunInTime_lengthBits}
\difficulty{3}
The function sending a Boolean word $u$ to the binary expansion of its length is
polynomial-time computable.
\end{theorem}

\begin{definition}[Constant-output map]
\lean{Complexity.Meyer.constMap}
\label{Complexity.Meyer.constMap}
\leanok
\uses{Complexity.transduce}
For a Boolean value $b$ and a bit list $u$, define a map that processes $u$ with a
transducer whose output at each input position is $b$.
\end{definition}

\begin{theorem}[Constant transduction]
\lean{Complexity.Meyer.constMap_eq}
\label{Complexity.Meyer.constMap_eq}
\leanok
\uses{Complexity.Meyer.constMap, Complexity.transduce}
\difficulty{4}
For every Boolean bit $b$ and finite Boolean word $u$, the one-pass transducer that
emits $b$ on each input symbol maps $u$ to the word consisting of $|u|$ copies of $b$.
\end{theorem}

\begin{theorem}[Polynomial-time computability of constant output]
\lean{Complexity.Meyer.pt_constMap}
\label{Complexity.Meyer.pt_constMap}
\leanok
\uses{Complexity.Meyer.constMap, Complexity.PolyTimeComputable, Complexity.polyTimeComputable_transduce}
\difficulty{1}
For each bit $b$, the function sending a Boolean word $u$ to the word consisting of
$|u|$ copies of $b$ is polynomial-time computable.
\end{theorem}

\begin{definition}[Partial Boolean state transition]
\lean{Complexity.Meyer.fTδ}
\leanok
For a state in the nested option type and an input bit, define a partial transition as
follows: an absent state remains absent; a state of the form “present, then absent”
becomes the input bit stored inside two present option constructors; and a state
containing a bit produces an absent result when that bit equals the input, and otherwise
produces an outer present constructor containing an absent one.
\end{definition}

\begin{definition}[Boolean output transition]
\lean{Complexity.Meyer.fTo}
\label{Complexity.Meyer.fTo}
\leanok
For a state in the nested option type and an input bit, define an output bit as follows:
if the state is absent or is “present, then absent,” the output is the input bit; if the
state contains a bit, the output is true exactly when that bit equals the input.
\end{definition}

\begin{definition}[Boolean transducer on words]
\lean{Complexity.Meyer.fT}
\label{Complexity.Meyer.fT}
\leanok
\uses{Complexity.Meyer.fTo, Complexity.Meyer.fTδ, Complexity.transduce}
For a Boolean word $u$, define its output by running the transducer with the partial
state transition and output rule described by the given transition functions, starting
in the state “present, then absent,” on the input word $u$.
\end{definition}

\begin{theorem}[Copying the input word]
\lean{Complexity.Meyer.transduce_fT_copy}
\label{Complexity.Meyer.transduce_fT_copy}
\leanok
\uses{Complexity.Meyer.fTo, Complexity.Meyer.fTδ, Complexity.transduce}
\difficulty{4}
For every finite Boolean word $u$, the transducer with transition rule $fTδ$ and
emission rule $fTo$, started in state $\mathrm{none}$, outputs $u$ itself.
\end{theorem}

\begin{theorem}[Transduction of a paired word]
\lean{Complexity.Meyer.fT_pairEncode}
\label{Complexity.Meyer.fT_pairEncode}
\leanok
\uses{Complexity.Meyer.fT, Complexity.Meyer.fTo, Complexity.Meyer.fTδ, Complexity.Meyer.fcode, Complexity.Meyer.transduce_fT_copy, Complexity.transduce, Turing.pairEncode}
\difficulty{4}
For finite Boolean words $w$ and $r$, let $F$ be the one-pass bit transduction specified
by the transition and output rules that define $fT$, and let $\langle w,r\rangle$ be the
encoding formed by doubling each bit of $w$, appending the separator $01$, and then
appending $r$. Applying $F$ to this encoding yields the word obtained by replacing each
bit $b$ of $w$ by $b1$, followed by $00$ and then $r$.
\end{theorem}

\begin{theorem}[Polynomial-time computability of a finite-state transduction]
\lean{Complexity.Meyer.pt_fT}
\label{Complexity.Meyer.pt_fT}
\leanok
\uses{Complexity.Meyer.fT, Complexity.PolyTimeComputable, Complexity.polyTimeComputable_transduce}
\difficulty{1}
The function on Boolean words computed by the following finite-state transducer is
polynomial-time computable. Its state records either no pending bit, a pending first bit
of a pair, or a state in which it emits each input bit unchanged; initially it has no
pending bit, and it emits each first bit and then, for a matching second bit, the bit
indicating whether the two bits agree. If the pair disagrees, subsequent bits are
emitted unchanged.
\end{theorem}

\begin{theorem}[Polynomial-time closure under query formation]
\lean{Complexity.Meyer.pt_query}
\label{Complexity.Meyer.pt_query}
\leanok
\uses{Complexity.Meyer.fT, Complexity.Meyer.fT_pairEncode, Complexity.Meyer.pt_fT, Complexity.Meyer.query, Complexity.PolyTimeComputable, Complexity.PolyTimeComputable.comp, Complexity.PolyTimeComputable.pairEncode, Complexity.polyTimeComputable_prepend, Turing.pairEncode}
\difficulty{4}
Fix a finite bit vector $b\in\{0,1\}^{n_b}$, and let $T,A,X:\{0,1\}^*\to\{0,1\}^*$ be
polynomial-time computable. For a bit string $z$, form the concatenation of $b$, the
encoding of $T(z)$ obtained by doubling each bit, the separator $00$, the similarly
encoded $A(z)$, another separator $00$, and $X(z)$. The resulting function of $z$ is
polynomial-time computable.
\end{theorem}

\begin{definition}[Polynomial growth bound]
\lean{Complexity.Meyer.wd}
\label{Complexity.Meyer.wd}
\leanok
For natural numbers $C$, $c$, and $n$, define $W(C,c,n)=C(n+1)^c$.
\end{definition}

\begin{definition}[Four-symbol tag type]
\lean{Complexity.Meyer.TW}
\label{Complexity.Meyer.TW}
\leanok
Let $TW$ be a type with exactly four elements, denoted $t_c$, $t_n$, $t_0$, and $t_1$.
\end{definition}

\begin{definition}[Six-symbol tag type]
\lean{Complexity.Meyer.AW}
\label{Complexity.Meyer.AW}
\leanok
Let $AW$ be a type with exactly six elements, denoted $a_c$, $a_n$, $a_0$, $a_1$, $a_j$,
and $a_{j1}$.
\end{definition}

\begin{definition}[Encoded length and padding word]
\lean{Complexity.Meyer.bjw}
\label{Complexity.Meyer.bjw}
\leanok
\uses{Complexity.Meyer.vj, Complexity.Meyer.vp}
For natural numbers $C_w,c_w$ and a Boolean word $z$, let $v_j(z)$ and $v_p(z)$ be the
natural numbers used by the given encoding. Define $b(z)$ to be the binary expansion of
$v_j(z)$, followed by a word of $v_p(z)$ zeros.
\end{definition}

\begin{definition}[Tag-dependent encoded word]
\lean{Complexity.Meyer.twd}
\label{Complexity.Meyer.twd}
\leanok
\uses{Complexity.Meyer.TW, Complexity.Meyer.vt, Complexity.Meyer.vx, Complexity.Meyer.wd, Turing.incFixed}
Let $C_w,c_w$ be natural numbers and $z$ a Boolean word. Write $x(z)=\pi_1(\pi_1(z))$
and $t(z)=\pi_1(\pi_2(z))$, using the total pair projections. For tag $t_c$, return
$t(z)$; for $t_n$, return its fixed-width binary successor, or the empty word if
increment overflows; for $t_0$ and $t_1$, return respectively a word of
$C_w(|x(z)|+1)^{c_w}$ zeros or ones.
\end{definition}

\begin{definition}[Outputs for alternating-word types]
\lean{Complexity.Meyer.awd}
\label{Complexity.Meyer.awd}
\leanok
\uses{Complexity.Meyer.AW, Complexity.Meyer.bjw, Complexity.Meyer.va, Complexity.Meyer.vx, Complexity.Meyer.wd, Turing.incFixed}
Fix natural numbers $C_w$ and $c_w$. For a type in
$\{\mathrm{ac},\mathrm{an},\mathrm{a0},\mathrm{a1},\mathrm{aj},\mathrm{aj1}\}$ and a
Boolean word $z$, define its output word as follows: for $\mathrm{ac}$ use the parsed
second component of the outer second projection of $z$; for $\mathrm{an}$ use its
fixed-width binary successor, taking the empty word on overflow; for $\mathrm{a0}$ use
the all-zero word of length $C_w(|x|+1)^{c_w}$, where $x$ is the first component of the
first projection of $z$; for $\mathrm{a1}$ use the successor of that all-zero word,
again taking the empty word on overflow. For $\mathrm{aj}$ use the concatenation of the
binary representation of the length of the parsed first component of the outer second
projection and an all-zero word whose length is the parsed second component's length;
for $\mathrm{aj1}$ use its fixed-width successor, taking the empty word on overflow.
Every parsed component here is obtained by the total projections, which return the empty
word when parsing fails.
\end{definition}

\begin{theorem}[Polynomial-time bounded replication]
\lean{Complexity.Meyer.pt_rep}
\label{Complexity.Meyer.pt_rep}
\leanok
\uses{Complexity.Meyer.pt_constMap, Complexity.Meyer.pt_vx, Complexity.Meyer.vx, Complexity.Meyer.wd, Complexity.PolyHierarchy.unaryPT_poly, Complexity.PolyTimeComputable, Complexity.PolyTimeComputable.comp}
\difficulty{3}
Let $C_w,c\in\mathbb{N}$ and fix $b\in\{0,1\}$. For a bit string $z$, apply the total
first projection twice to obtain a bit string $x$ (each projection returns the empty
word when its input does not parse as a pair). Then output the constant word consisting
of $b$ repeated $C_w(|x|+1)^c$ times. This function of $z$ is polynomial-time
computable.
\end{theorem}

\begin{theorem}[Polynomial-time computability of the length encoding]
\lean{Complexity.Meyer.pt_bjw}
\label{Complexity.Meyer.pt_bjw}
\leanok
\uses{Complexity.Meyer.bjw, Complexity.Meyer.pt_constMap, Complexity.Meyer.pt_lengthBits, Complexity.Meyer.pt_vj, Complexity.Meyer.pt_vp, Complexity.PolyTimeComputable, Complexity.PolyTimeComputable.append, Complexity.PolyTimeComputable.comp, Complexity.polyTimeComputable_unary}
\difficulty{3}
The function sending a Boolean word $z$ to the binary expansion of the length of the
second component extracted by the total second projection of the encoded second and
third components of $z$, followed by a string of zeros whose length is the length of the
total fourth projection of $z$, is polynomial-time computable.
\end{theorem}

\begin{theorem}[Polynomial-time computability of four transformations]
\lean{Complexity.Meyer.pt_twd}
\label{Complexity.Meyer.pt_twd}
\leanok
\uses{Complexity.Meyer.TW, Complexity.Meyer.pt_inc, Complexity.Meyer.pt_rep, Complexity.Meyer.pt_vt, Complexity.Meyer.twd, Complexity.PolyTimeComputable, Complexity.PolyTimeComputable.comp}
\difficulty{4}
Let $C_w,c\in\mathbb{N}$. For every choice among four transformations of bit strings,
the resulting function is polynomial-time computable: return the second component of the
twice nested pair parsing; return its fixed-width binary successor, or the empty word if
that increment overflows or the parsing fails; or return a word of $C_w(|x|+1)^c$ zeros
or ones, where $x$ is obtained by applying the total first projection twice (with failed
parses yielding the empty word).
\end{theorem}

\begin{theorem}[Polynomial-time computability of the AW operations]
\lean{Complexity.Meyer.pt_awd}
\label{Complexity.Meyer.pt_awd}
\leanok
\uses{Complexity.Meyer.AW, Complexity.Meyer.awd, Complexity.Meyer.pt_bjw, Complexity.Meyer.pt_inc, Complexity.Meyer.pt_rep, Complexity.Meyer.pt_va, Complexity.PolyTimeComputable, Complexity.PolyTimeComputable.comp}
\difficulty{4}
For all natural-number parameters $C_w,c_w$ and every operation $A$ in the six-element
set consisting of copying the second component of a decoded input, incrementing that
component, producing an all-zero word of length $C_w(|x|+1)^{c_w}$, incrementing that
all-zero word, producing the word of binary digits of the length of the fourth decoded
component followed by zeros whose number is that component's length, or incrementing
this last word, the resulting function on Boolean words is polynomial-time computable.
In each increment case, an overflow is replaced by the empty word.
\end{theorem}

\begin{definition}[Index description word]
\lean{Complexity.Meyer.idxDesc}
\label{Complexity.Meyer.idxDesc}
\leanok
For natural numbers $n$ and $j$, let the index description be the Boolean word
consisting of $n$ true bits, followed by two false bits, then $j$ true bits, and finally
one false bit.
\end{definition}
